\section{Answer 1: what the entry of a power of \texorpdfstring{$A$}{A} counts}\label{sec:answer1}
% claims used here: C3 C6 C7 C8 C9 C10 C11 m1-2 m1-3 m1-5 m1-6 m1-7 m1-8 m1-9 m3-11

Question~1 asks what $(A^k)_{ij}$ represents. The answer is a count of walks
(\Cref{def:walk}), and the proof needs only the definition of the matrix
product, so that product comes first.

\subsection{Matrix products and powers}\label{sec:product}

\begin{plainwords}
To multiply two tables, take a row from the first and a column from the second.
Multiply the numbers in pairs. Then add the results. Do this for every row and
every column to build the new table. Multiply the road table by itself again
and again, and you get its powers.
\end{plainwords}

\begin{definition}[Matrix product and power]\label{def:product}
Let $X$ and $Y$ be $N\times N$ matrices. Their \emph{product} $XY$ is the $N\times N$ matrix with
\begin{equation}\label{eq:product}
(XY)_{ij}=\sum_{r=1}^{N}X_{ir}\,Y_{rj}.
\end{equation}
The \emph{powers} of a square matrix $A$ are $A^1=A$ and $A^k=A^{k-1}A$ for $k\ge2$.
\end{definition}

\begin{remark}\label{rem:product}
Formula~\eqref{eq:product} multiplies row $i$ of $X$ with column $j$ of $Y$ entry by entry and adds. A power is the repeated product of the whole matrix, not a power of each entry. This paper uses exponents $k\ge1$ only.
\end{remark}

\begin{example}\label{ex:product}
For the board matrix, $(A^2)_{22}=A_{21}A_{12}+A_{22}A_{22}+A_{23}A_{32}=1\cdot1+0\cdot0+1\cdot1=2$ (by script, \cref{app:scripts}).
\end{example}

\begin{nonexample}\label{nex:product}
$A^2$ is not the matrix of squared entries. Squaring each entry of the board matrix gives back $A$, whose entry $(1,2)$ is $1$, while $(A^2)_{12}=0$.
\end{nonexample}

\subsection{The walk-counting theorem}\label{sec:walkcount}

\begin{plainwords}
Take the table of roads and multiply it by itself. Each number in the new table
counts trips of two roads. The trips run from the row house to the column house.
Multiply again and you count trips of three roads. The count includes trips that
turn back. So the table counts trips, not routes.
\end{plainwords}

\begin{theorem}[Walk counting]\label{thm:walkcount}
Let $A$ be the adjacency matrix of a simple undirected graph $G$ with vertices $n_1,\dots,n_N$, and let $k\ge1$ be an integer. Then $(A^k)_{ij}$ is the number of walks of length $k$ from $n_i$ to $n_j$, for all $i,j$.
\end{theorem}

The statement is Theorem~1.1.3 of the University of Galway notes, for exactly this setting~\cite[Theorem~1.1.3]{quinlan_matrices_graphs}. Grinberg, Lehman, Leighton and Meyer, and Bai state it too, in more general settings~\cite[Theorem~4.5.10, pp.~120--121]{grinberg2025graph}\cite[Corollary~9.3.3]{lehman_mcs_adjacency}\cite{bai2025adjacency}. The textbook does not state the theorem in the pages checked \kindtag{qualified}, so the proof below is written out in full. It follows the standard route: induction on $k$, splitting the walks by their second-last vertex~\cite{quinlan_matrices_graphs,grinberg2025graph}. The write-up is \kindtag{derived here}. One journal paper states a different version, with ``paths'' in place of walks; it is false, and \cref{sec:warning} shows why.

\begin{remark}\label{rem:walkcount}
The theorem allows $i=j$, which counts closed walks, and it counts every walk once, with repeated vertices and edges allowed. It forbids reading the number as a count of paths (\cref{prop:warning}). It says nothing about walks of other lengths: $(A^2)_{12}=0$ on the board graph does not say that $n_1$ and $n_2$ are unconnected.
\end{remark}

\begin{example}\label{ex:walkcount}
On the board graph $B$, $(A^2)_{13}=1$: the only walk of length $2$ from $n_1$ to $n_3$ is $1,2,3$. And $(A^3)_{12}=2$: the walks of length $3$ from $n_1$ to $n_2$ are $1,2,1,2$ and $1,2,3,2$. All entries are written out in \cref{sec:example}.
\end{example}

\begin{nonexample}\label{nex:walkcount}
$(A^2)_{12}=0$ on $B$ does not mean that no walk joins $n_1$ and $n_2$. It means only that no walk of length \emph{exactly} $2$ does; walks of length $1$ and of length $3$ exist.
\end{nonexample}

\begin{proof}
We use induction on $k$. Write $W_k(i,j)$ for the number of walks of length $k$ from $n_i$ to $n_j$. We prove $W_k(i,j)=(A^k)_{ij}$ for all $i,j$, for each $k\ge1$.
In a simple graph a walk is determined by its list of vertices (\cref{rem:walk}). So $W_k(i,j)$ is the number of lists $(w_0,\dots,w_k)$ of vertices with $w_0=n_i$, $w_k=n_j$, and $w_{t-1}$ adjacent to $w_t$ for each $t$ from $1$ to $k$.

\emph{Base case, $k=1$.} A walk of length $1$ from $n_i$ to $n_j$ is a list $(n_i,n_j)$ with $n_i$ and $n_j$ adjacent. There is exactly one such list if they are adjacent and none otherwise. So $W_1(i,j)=A_{ij}=(A^1)_{ij}$.

\emph{Induction hypothesis.} Let $k\ge2$ and suppose $W_{k-1}(i,j)=(A^{k-1})_{ij}$ for all $i,j$.

\emph{Induction step.} Fix $i$ and $j$. Every walk $(w_0,\dots,w_k)$ of length $k$ from $n_i$ to $n_j$ has a second-last vertex $w_{k-1}$. Because $k\ge2$, deleting the last vertex leaves $w_0,\dots,w_{k-1}$, which still has length $k-1\ge1$, so the hypothesis applies to it. (For $k=1$ nothing would remain to count; this is why $k=1$ is the separate base case.) Let $w_{k-1}=n_r$, where $r$ is one of $1,\dots,N$. Sort the walks into $N$ classes by the value of $r$. The classes are disjoint, since each walk has exactly one second-last vertex, and together they contain every walk. So
\[
W_k(i,j)=\sum_{r=1}^{N}\bigl(\text{number of walks of length $k$ from $n_i$ to $n_j$ with second-last vertex $n_r$}\bigr).
\]
We claim that the number in class $r$ is $W_{k-1}(i,r)\,A_{rj}$.
\begin{itemize}
\item If $n_r$ and $n_j$ are not adjacent, then $A_{rj}=0$. No walk can have $n_r$ followed by $n_j$, so class $r$ is empty, and its count $0$ equals $W_{k-1}(i,r)\cdot0$.
\item If $n_r$ and $n_j$ are adjacent, then $A_{rj}=1$. Delete the last vertex of a walk in class $r$. What remains is a list $w_0,\dots,w_{k-1}$ with $w_0=n_i$, $w_{k-1}=n_r$ and consecutive vertices adjacent, so it is a walk of length $k-1$ from $n_i$ to $n_r$. Conversely, take any walk of length $k-1$ from $n_i$ to $n_r$ and append $n_j$. The new pair $n_r,n_j$ is adjacent, so the result is a walk of length $k$ from $n_i$ to $n_j$ with second-last vertex $n_r$. Deleting and appending undo each other, so they match the walks in class $r$ one to one with the walks of length $k-1$ from $n_i$ to $n_r$. The count is $W_{k-1}(i,r)=W_{k-1}(i,r)\cdot1$.
\end{itemize}
Therefore $W_k(i,j)=\sum_{r}W_{k-1}(i,r)\,A_{rj}$. By the induction hypothesis $W_{k-1}(i,r)=(A^{k-1})_{ir}$, so
\[
W_k(i,j)=\sum_{r=1}^{N}(A^{k-1})_{ir}\,A_{rj}=(A^{k-1}A)_{ij}=(A^k)_{ij},
\]
by formula~\eqref{eq:product} and the definition of powers. This proves the statement for $k$, and so, by induction, for every $k\ge1$.
\end{proof}

\paragraph{Why the argument is valid.}
Two facts carry it: the sum over the second-last vertex covers each walk exactly once, and the matching in the case $A_{rj}=1$ is one to one. Simplicity is used twice: to identify a walk with its vertex list, and so that $A_{rj}$, which is the number of edges between $n_r$ and $n_j$, is $0$ or $1$. Undirectedness is not used inside the induction; it only makes $A$ symmetric. \kindtag{interpretation} The proof is checked, not replaced, by a finite test. A \emph{case} is one choice of a graph, an exponent $k$ and a pair $i,j$; in each case a script counted the walks by brute force and compared the count with the entry $(A^k)_{ij}$. For every labelled simple graph on $2,3,4,5$ vertices ($2$, $8$, $64$, $1024$ graphs), with $k$ up to $5$ for $N\le4$ and up to $3$ for $N=5$, there are $2\cdot5\cdot4=40$, $8\cdot5\cdot9=360$, $64\cdot5\cdot16=5120$ and $1024\cdot3\cdot25=76800$ cases, and all are equal (\cref{app:crosschecks}). That test is a check, not a proof.

\subsection{Consequences of the theorem}\label{sec:consequences}

\begin{plainwords}
A count is never negative and always a whole number. The table of trips looks the
same when flipped over its diagonal. A positive number means that at least one
such trip exists. For trips of two roads, the diagonal counts the roads at a
house.
\end{plainwords}

\begin{corollary}\label{cor:walkcount}
For a simple undirected graph and $k\ge1$:
\begin{enumerate}
\item[(a)] every entry of $A^k$ is a non-negative integer, so a sum of such powers has a positive $(i,j)$ entry if and only if at least one summand does;
\item[(b)] $(A^k)_{ij}>0$ if and only if there is a walk of length exactly $k$ from $n_i$ to $n_j$;
\item[(c)] $A^k$ is symmetric;
\item[(d)] $(A^2)_{ii}=\deg(n_i)$, so the sum of the diagonal entries of $A^2$ equals $2|E|$.
\end{enumerate}
\kindtag{derived here}
\end{corollary}

\begin{proof}
Parts (a) and (b) follow because by \cref{thm:walkcount} each entry is a count. For (c), $(A^k)^{T}=(A^{T})^k=A^k$, where the \emph{transpose} $X^{T}$ of a matrix $X$ is the matrix with its rows and columns exchanged, and $A^{T}=A$ because $A$ is symmetric (\cref{prop:Aproperties}); alternatively, reversing a walk matches the walks from $n_i$ to $n_j$ with the walks from $n_j$ to $n_i$. For (d), the walks of length $2$ from $n_i$ to $n_i$ are the lists $n_i,n_r,n_i$ with $n_r$ a neighbour of $n_i$, one for each neighbour, and the number of neighbours is $\deg(n_i)$ (\cref{rem:degree}). Summing over $i$ and using part~(b) of \cref{prop:rowsums} gives $2|E|$. The University of Galway notes state (d) in a remark after their theorem~\cite{quinlan_matrices_graphs}.
\end{proof}

On the board graph the diagonal of $A^2$ is $(1,2,1)$, the degrees, and its sum is $4=2|E|$ (script).

\subsection{The warning: walks are not paths}\label{sec:warning}

\begin{plainwords}
A trip may turn back and pass the same house twice. A route never passes a
house twice. The table counts trips, so it also counts the trips that turn back.
A house with at least one neighbour always has a two-road trip that goes out and
comes home. That trip is no route.
\end{plainwords}

\begin{proposition}[The entry is not a count of paths]\label{prop:warning}
Let $G$ be a simple undirected graph.
\begin{enumerate}
\item[(a)] If $n_i$ has at least one neighbour, then $(A^2)_{ii}=\deg(n_i)>0$, yet there is no path, and no trail, of length $2$ from $n_i$ to $n_i$.
\item[(b)] On the board graph $B$, $(A^3)_{12}=2$, although the only path from $n_1$ to $n_2$ has length $1$. More generally, on $B$ (and not for every graph) $A^{k+2}=2A^k$ for $k\ge1$, so the entries grow with $k$ while the number of paths does not.
\end{enumerate}
\kindtag{derived here}
\end{proposition}

\begin{proof}
(a) By \cref{cor:walkcount}(d), $(A^2)_{ii}=\deg(n_i)$, which is positive. The only walks of length $2$ from $n_i$ to $n_i$ are $n_i,n_r,n_i$. They repeat the vertex $n_i$, so they are not paths (\cref{def:path}); a path from a vertex to itself has pairwise distinct vertices, so its length is $0$. They repeat the edge $n_in_r$, so they are not trails either (\cref{def:trail}). The textbook's own word ``path'' (Definition~1.2.12) means a trail, so the entry is not a count of the textbook's paths either.
(b) On $B$ the walks of length $3$ from $n_1$ to $n_2$ are $1,2,1,2$ and $1,2,3,2$. Each repeats a vertex and an edge, so neither is a path or a trail. The only path from $n_1$ to $n_2$ is the single edge $1,2$. For the growth, $A^3=2A$ (script) and multiplying by $A^{k-1}$ gives $A^{k+2}=2A^k$.
\end{proof}

The sources say the same in words: Bai remarks that the entries count walks and not paths~\cite{bai2025adjacency}, and Grinberg's Example~2.9.2 shows a walk that is not a path~\cite[p.~37]{grinberg2025graph} . \Cref{tab:walksvspaths} gives the counts from the script for $k=1,\dots,5$: the walk counts $1,0,2,0,4$ against the path counts $1,0,0,0,0$.

\begin{table}[ht]
\centering
\caption{Walks of length $k$ from $n_1$ to $n_2$ on the board graph $B$, and how many of them are paths and trails (by script). The walk counts equal the entries $(A^k)_{12}$ of \cref{sec:example}.}
\label{tab:walksvspaths}
\begin{tabular}{@{}cccc@{}}
\toprule
$k$ & walks & of which paths & of which trails \\
\midrule
\TabWalksVsPathsBody
\bottomrule
\end{tabular}
\end{table}

Paramadevan and Sotheeswaran's Theorem~1 reads: ``the number of different paths of length $l$ between the $i$-th and $j$-th vertices of $G$ is equal to the $(i,j)$-th entry of $A_G^l$''~\cite{paramadevan2021adjacency}. As written this is false for any definition of path that forbids repeated vertices or repeated edges, by \cref{prop:warning}(a); their own property~7 says walks~\cite{paramadevan2021adjacency}. So the sources disagree, and the paper follows the walk version, which four sources state and which was proved above.

What a positive entry \emph{does} guarantee is a walk of length $k$, and then, by \cref{lem:walkpath}, a path of length at most $k$. This is why the warning does not harm the connectivity test of \cref{sec:answer2}; but the value of the entry is a count of walks, not of paths. \kindtag{derived here}
